\documentclass[12pt]{article}
\usepackage[margin=1.2in]{geometry}
\usepackage{fontspec}
\usepackage{titling}
\usepackage{verse}
\usepackage{xcolor}

\setmainfont{EB Garamond}

\definecolor{desert}{RGB}{26,14,8}
\definecolor{cactus}{RGB}{107,142,35}
\definecolor{failure}{RGB}{204,85,64}
\definecolor{gold}{RGB}{212,168,48}
\definecolor{text}{RGB}{232,224,216}

\pagecolor{desert}
\color{text}

\title{\textcolor{gold}{The Parking Orbit}}
\author{Monologues from Dead Spacecraft\\[0.3em]\small\textit{After Edgar Lee Masters}}
\date{Mission 1.26 --- Ranger 1 + \textit{Opuntia fragilis}\\April 5, 2026}

\begin{document}
\maketitle

\begin{center}
\textcolor{cactus}{\textit{``Where are Elmer, Herman, Bert, Tom and Charley,\\
The weak of will, the strong of arm, the clown, the boozer, the fighter?''}}\\
\small --- Edgar Lee Masters, \textit{Spoon River Anthology}, ``The Hill''
\end{center}

\vspace{1.5em}

\section*{\textcolor{failure}{Ranger 1}}

\begin{verse}
I had the fuel. I had the engine.\\
I had the trajectory computed to seven decimals\\
and the burn sequence programmed into the timer\\
and the oxidizer loaded to the milligram.\\
What I did not have was gravity.\\
In the parking orbit, 160 kilometers above the town,\\
the propellant floated like a man\\
who has lost the direction of down.\\
The valves opened. The chamber waited.\\
The liquid drifted past the inlet\\
carrying bubbles of nitrogen pressurant\\
where there should have been hydrazine.\\
I circled. One hundred and fifteen times.\\
Each pass the air a little thicker,\\
each pass the ground a little closer.\\
On the seventh day I came home the hard way ---\\
three hundred and six kilograms of capability\\
resolved into a streak of light\\
that no one on the ground mistook for a star.
\end{verse}

\vspace{1.5em}

\section*{\textcolor{failure}{Pioneer 0}}

\begin{verse}
Seventy-seven seconds.\\
That is what I was given.\\
The bearing in the turbopump seized\\
and the fuel line starved\\
and the thrust chamber coughed once\\
and the vehicle broke apart at Max Q\\
and I fell into the Atlantic\\
in pieces too small to recover.\\
They found the bearing.\\
They wrote the specification.\\
They fixed it for the next one.\\
I was the specification.
\end{verse}

\vspace{1.5em}

\section*{\textcolor{failure}{Ranger 6}}

\begin{verse}
I did everything right.\\
The Atlas burned clean. The Agena restarted.\\
The trajectory was perfect --- I hit the Moon\\
within two kilometers of the target,\\
which is the orbital mechanics equivalent\\
of threading a needle from across a room\\
while the room is spinning.\\
My spacecraft bus worked. The solar panels,\\
the attitude control, the telemetry ---\\
all nominal for three days of cislunar cruise.\\
I arrived at the Moon precisely on schedule.\\
The six television cameras did not turn on.\\
A short circuit during the Atlas staging\\
had fused the camera power bus\\
seventy-two hours before I needed it.\\
I hit the Moon blind.\\
The engineers found the short.\\
They redesigned the power routing.\\
Ranger 7 carried the same cameras\\
and returned four thousand three hundred and sixteen photographs\\
of the surface I saw only with my hull.
\end{verse}

\vspace{1.5em}

\section*{\textcolor{failure}{The Agena A Engine}}
\small\textit{Bell 8081, Serial 6001}

\begin{verse}
They blame me.\\
The engine that would not restart.\\
But I was not built to restart.\\
I was built to ignite once,\\
in the presence of gravity,\\
with liquid propellant settled at the bottom\\
of tanks that assumed a bottom existed.\\
They asked me to ignite in a place\\
where bottom is a meaningless word.\\
Where the fuel floats like thoughts\\
in the mind of a man who cannot decide.\\
I opened my valves. I waited for propellant.\\
What arrived was a mixture of hope and nitrogen.\\
I could not burn hope.
\end{verse}

\vspace{2em}

\begin{center}
\textcolor{cactus}{\textit{They speak from their orbits. Some are still circling.\\
Some burned up. Some are on the Moon.\\
Each one was the next try.}}
\end{center}

\vspace{1em}

\begin{center}
\small\textcolor{gold}{NASA Mission Series v1.26 --- Research by Miles Tiberius Foxglove --- Mukilteo, WA}
\end{center}

\end{document}
